\documentclass[12pt]{article}
\usepackage[spanish]{babel}
\usepackage[utf8]{inputenc}
\usepackage[T1]{fontenc}
\usepackage{geometry}
\geometry{letterpaper,margin=1in}
\usepackage{amsmath,amssymb,mathtools}
\usepackage{siunitx}
\usepackage{bm}
\usepackage{physics}
\usepackage{graphicx}
\usepackage{enumitem}
\setlist[itemize]{topsep=2pt,itemsep=2pt}
\setlist[enumerate]{topsep=2pt,itemsep=2pt}

\title{\Large Problemas (Sección 8) \\[2mm]
\small Capítulo I --- Ecuaciones de Lagrange}
\author{Transcripción a \LaTeX}
\date{\today}

\begin{document}
\maketitle

\section*{8.\ Problemas}

\paragraph{1.1}\textbf{(Partícula en coordenadas polares)}\;
Una partícula de masa $m$ se mueve en un plano bajo la influencia de una fuerza que, en coordenadas polares $(r,\theta)$, está dada por
\begin{equation}
\label{eq:F11}
\mathbf F \;=\; \frac{1}{r^{2}}\!\left( 1 - \frac{\dot r^{2} - 2 r \ddot r}{c^{2}} \right)\hat{\mathbf r}.
\end{equation}
Determinar el \emph{potencial generalizado} $U(r,\dot r)$ que proporciona esta fuerza y la Lagrangiana que determina la evolución de la partícula en el plano.

\paragraph{1.2}\textbf{(Lagrangiana cuadrática en $x,y$)}\;
Una Lagrangiana para un sistema físico está dada por
\begin{equation}
\label{eq:L12}
L \;=\; \frac{m}{2}\!\left(a\,\dot x^{2} + 2b\,\dot x \dot y + c\,\dot y^{2}\right)
- \frac{K}{2}\!\left(a\,x^{2} + 2b\,x y + c\,y^{2}\right).
\end{equation}
donde $a,b,c$ son constantes reales arbitrarias tales que $b^{2}-ac\neq 0$.
\begin{itemize}
  \item Obtener las ecuaciones de movimiento.
  \item Estudiar con detalle los casos particulares $a=0=c$ y $b=0,\; c=-a$.
  \item Físicamente, ¿qué describe la Lagrangiana \eqref{eq:L12}?
\end{itemize}

\paragraph{1.3}\textbf{(Lagrangiana no usual en 1D)}\;
Una partícula de masa $m$ se mueve en una dimensión; su Lagrangiana está dada por
\emph{(véase la figura siguiente, recorte fiel del enunciado en el PDF)}.
\begin{center}
\includegraphics[width=.9\linewidth]{problem_1_3.png}
\end{center}
\noindent
\textit{Pedir:} Obtener la ecuación de movimiento y describir la naturaleza física del sistema.

\paragraph{1.4}\textbf{(Invariancia de norma en electromagnetismo)}\;
Muestre que el electromagnetismo clásico es invariante ante la transformación de \emph{calibre} (gauge) de los potenciales
\begin{equation}
\label{eq:gauge}
\mathbf A \to \mathbf A + \nabla \psi(\mathbf r,t),
\qquad
\phi \to \phi - \frac{1}{c}\,\frac{\partial \psi(\mathbf r,t)}{\partial t},
\end{equation}
donde $\psi(\mathbf r,t)$ es una función diferenciable.
Determinar el cambio de la Lagrangiana de una partícula que se mueve en un campo electromagnético ante una transformación de norma. 
¿El movimiento de la partícula se ve afectado?

\paragraph{1.5}\textbf{(Potencial generalizado $U(r,\mathbf v)=V(r)+\mathbf a\cdot\mathbf L$)}\;
Una partícula de masa $m$ se mueve en el espacio bajo la influencia de una fuerza derivable del potencial generalizado
\begin{equation}
U(\mathbf r,\mathbf v) \;=\; V(r) + \mathbf a \cdot \mathbf L,
\end{equation}
donde $\mathbf r$ es el vector de posición de la partícula (respecto a un sistema inercial), $\mathbf L$ es su momento angular y $\mathbf a$ es un vector fijo en el espacio.
\begin{enumerate}[label=(\alph*)]
\item Determinar las componentes de la fuerza en coordenadas Cartesianas y en coordenadas esféricas.
\item Obtener las ecuaciones de movimiento en coordenadas esféricas.
\end{enumerate}

\bigskip
\hrule
\bigskip

\paragraph{1.6}\textbf{(Péndulo esférico)}\;
Obtener las ecuaciones de Lagrange de un \emph{péndulo esférico}; es decir, una partícula de masa $m$ suspendida por una varilla de masa despreciable sujeta al campo gravitacional de la Tierra.

\paragraph{1.7}\textbf{(Dos masas unidas por barra rígida)}\;
Dos partículas puntuales de masa $m$ están unidas por una barra rígida de masa despreciable de longitud $l$, cuyo centro está obligado a moverse en un círculo de radio $a$. El sistema se mueve en el espacio bajo la acción del campo gravitacional de la Tierra.
Obtener la Lagrangiana del sistema y las ecuaciones de movimiento.

\paragraph{1.8}\textbf{(Derivada total añadida al Lagrangiano)}\;
Si $L$ es un Lagrangiano para un sistema de $n$ grados de libertad que satisface las ecuaciones de Lagrange, demuestre por sustitución directa que
\begin{equation}
L' \;=\; L + \frac{d}{dt}\,F(q_i,t)
\end{equation}
también satisface las ecuaciones de Lagrange, donde $F$ es una función arbitraria y diferenciable de sus argumentos.

\paragraph{1.9}\textbf{(Transformación de punto)}\;
Sea $\{q_i\}_{i=1}^n$ un conjunto de coordenadas generalizadas para un sistema de $n$ grados de libertad, con Lagrangiano $L$.
Supongamos que transformamos a otro conjunto $\{s_j\}_{j=1}^n$ mediante la transformación de punto
\begin{equation}
q_i \;=\; q_i(s_1,\ldots,s_n,t).
\end{equation}
Demuestre que, si la función Lagrangiana se expresa como $\tilde L(s,\dot s,t)=L(q(s,t),\dot q(s,\dot s,t),t)$ a través de la ecuación de transformación, entonces $\tilde L$ satisface las ecuaciones de Lagrange con respecto a las coordenadas $s_j$.
En otras palabras, la forma de las ecuaciones de Lagrange es invariante bajo una transformación de punto.

\paragraph{1.10}\textbf{(Dos ruedas con rodadura sin deslizamiento)}\;
Dos ruedas de radio $a$ están montadas en los extremos de un eje común de longitud $b$ de manera que las ruedas giran independientemente. El sistema rueda sin resbalar sobre un plano. 
Muestre que hay dos ecuaciones \emph{no holónomas} de constricción dadas por
\begin{align}
\cos\theta\,dx + \sin\theta\,dy &= 0, \\
\sin\theta\,dx - \cos\theta\,dy &= a\,(d\varphi + d\varphi'),
\end{align}
donde $\theta$ es el ángulo entre el eje común de longitud $b$ y el eje $x$ positivo; $\varphi$ y $\varphi'$ son los ángulos medidos alrededor del eje común $b$ que etiquetan los puntos de las ruedas; y $(x,y)$ son las coordenadas del punto medio del eje común entre las dos ruedas.
\par\noindent
Además, muestre que existe una ecuación \emph{holónoma} de constricción dada por
\begin{equation}
\theta \;=\; C - \frac{a}{b}\,(\varphi - \varphi'),
\end{equation}
donde $C$ es una constante.

\end{document}
